% Generated draft: requires mathematical and source review.
\documentclass[11pt]{article}
\usepackage[margin=25mm]{geometry}
\usepackage{fontspec}
\setmainfont{Latin Modern Roman}
\usepackage{amsmath,amssymb,mathtools,unicode-math}
\setmathfont{Latin Modern Math}
\usepackage{enumitem,needspace,xurl,hyperref}
\setlength{\parindent}{0pt}
\setlength{\parskip}{5pt}
\setlength{\emergencystretch}{3em}
\newcommand{\N}{\mathbb{N}}
\newcommand{\Z}{\mathbb{Z}}
\newcommand{\Q}{\mathbb{Q}}
\newcommand{\R}{\mathbb{R}}
\newcommand{\C}{\mathbb{C}}
\newcommand{\F}{\mathbb{F}}
\newcommand{\A}{\mathbb{A}}
\newcommand{\PP}{\mathbb{P}}
\newcommand{\E}{\mathbb{E}}
\newcommand{\eps}{\varepsilon}
\newcommand{\dd}{\,\mathrm{d}}
\begin{document}
\section*{UnsolvedMath ID: 20000236}
\textbf{Source problem code:} catalogue entry
\subsection{Definitions and mathematical statement}
Use the supplied contextual reconstruction: let \(E=H^0(\PP^1,\mathcal O(d))\), \(n=d+1\), \(k=r+1\), and \(G=\operatorname{Gr}(k,E)\), of dimension \(N=k(n-k)\). Let \(F_\bullet(x)\) be the complete osculating flag defined by vanishing orders at \(x\in\PP^1\). For pairwise distinct marked points \(x_1,\ldots,x_s\) and partitions \(\lambda_i\subseteq k\times(n-k)\), form the scheme-theoretic intersection
\[
S(\lambda_\bullet;x_\bullet)=\bigcap_{i=1}^s
\Omega_{\lambda_i}(F_\bullet(x_i)),\qquad
\sum_i|\lambda_i|=N-1.
\]
The supplied note attributes proper intersection to Eisenbud--Harris, making nonempty \(S\) a projective Cohen--Macaulay curve. The question asks whether \(S\) is smooth over \(\C\), including the absence of nonreal singular points. The intended restriction to real versus arbitrary distinct complex \(x_i\) is not completely recoverable: the contextual verification ends mid-sentence after ``All \(x_i\) are.''

The source proposes tautological bundles on \(M_{0,s}\) to describe the singular-fiber locus, possibly showing it avoids configurations with all \(x_i\) real. It reports a check for \(\operatorname{Gr}(2,5)\) with five single-box conditions, and suggests longer partitions or \(\operatorname{Gr}(2,6)\) next. Its sum-of-squares example uses terms \((x_i-x_j)^2\), suggesting collisions at the boundary. Such a real sum-of-squares exclusion alone does not establish absence of complex singular points. Its final ``B. Osserman suggests'' has no following suggestion in this record.

The attached aid reports two partial smoothness results: every one-dimensional osculating intersection with at most two nonempty conditions is a normal Richardson curve, hence complex-smooth without a reality assumption; and every real osculating Schubert curve of Pl\"ucker degree at most two is complex-smooth. Its explicit longer-partition examples on \(\operatorname{Gr}(2,6)\) are \((2,2)\) with three boxes and \((3,2)\) with two boxes, each of total codimension seven.

For the all-box case on \(\operatorname{Gr}(2,n)\), there are \(s=2n-5\) conditions. In a local coordinate at \(x_i\), let
\[
B_i(f,g)=f(x_i)g'(x_i)-f'(x_i)g(x_i)
\]
be the alternating Wronskian evaluation form, understood up to a nonzero fiber trivialization. The aid reports the exact singularity criterion for \(U\in S\):
\[
U\text{ singular in }S
\quad\Longleftrightarrow\quad
\exists[a_1:\cdots:a_s]\in\PP^{s-1},\quad
U\subseteq\operatorname{rad}\!\left(\sum_i a_iB_i\right).
\]
Here \(\operatorname{rad}(B)=\{v\in E:B(v,w)=0\text{ for all }w\in E\}\). It reports gluing this local description via Wronskian evaluation line bundles to a global projective incidence and an equivalent Wronski-map transversality criterion. These descriptions reduce the hidden-singularity question; they do not prove general complex smoothness.
\subsection{Short English statement}
Determine whether one-dimensional osculating Schubert intersections are complex-smooth, with careful attention to whether their marked points are real. The attached aid reports two-condition and low-Pl\"ucker-degree cases and an exact Wronskian-kernel incidence for all-box curves, leaving the general smoothness question open.
\subsection{Sources}
\begin{itemize}
\item Local UnsolvedMath registry, permanent ID \(20000236\), exact public code \texttt{catalogue entry}; canonical location \texttt{aim-algebraic-geometry-notes.json}, notes[235].
\item AIM workshop ``Degenerations in algebraic geometry,'' extracted Problem 10.4 and Example 10.5; supplied URL \url{https://aimath.org/pastworkshops/degenalggeomproblems.pdf}. No proposer or year is supplied; the literature field is empty. The source PDF was not consulted in this drafting step.
\item The supplied verification note reconstructs the preceding-record notation and cites Eisenbud--Harris proper intersection, but is truncated at the reality assumption. The source's unfinished attribution is ``B. Osserman suggests.''
\item Attached research attempt 1, ``Low-degree smoothness and a Wronskian kernel criterion for osculating Schubert curves,'' without complete bibliographic references.
\end{itemize}
\end{document}
